\section{Giới thiệu}

\subsection{Hiện trạng ngành nông nghiệp và chăm sóc cây trồng}
Trong bối cảnh nông nghiệp hiện đại, việc duy trì môi trường sống lý tưởng cho cây trồng là yếu tố tiên quyết quyết định năng suất và chất lượng nông sản. Hiện nay, tại nhiều hộ gia đình và trang trại quy mô vừa và nhỏ, quy trình chăm sóc vườn vẫn chủ yếu dựa vào sức người và kinh nghiệm cảm tính. Việc tưới nước thường được thực hiện thủ công hoặc thông qua các hệ thống bán tự động thô sơ, thiếu sự hỗ trợ của các thiết bị đo lường chính xác. Điều này không chỉ gây khó khăn trong việc quản lý diện tích lớn mà còn làm mất đi tính chủ động trong sản xuất trước những biến đổi thất thường của thời tiết.

\subsection{Những khó khăn và thách thức tồn đọng}
Các phương pháp quản lý vườn truyền thống hiện đang đối mặt với nhiều rào cản lớn, ảnh hưởng trực tiếp đến hiệu quả kinh tế:
\begin{itemize}
    \item \textbf{Lãng phí tài nguyên:} Việc không xác định được nhu cầu thực tế của cây trồng dẫn đến tình trạng tưới quá mức gây lãng phí nước, hoặc tưới thiếu làm cây bị stress và chậm phát triển.
    \item \textbf{Phụ thuộc vào sức lao động:} Người trồng phải có mặt trực tiếp để vận hành hệ thống, gây bất tiện khi cần di chuyển xa hoặc trong điều kiện thời tiết khắc nghiệt.
    \item \textbf{Thiếu dữ liệu giám sát:} Không có sự theo dõi liên tục về độ ẩm đất và các chỉ số môi trường, khiến việc xử lý các tình huống phát sinh như khô hạn hay ngập úng trở nên thụ động.
    \item \textbf{Rủi ro vận hành:} Các lỗi kỹ thuật như hỏng hóc cảm biến, cạn nguồn nước thường không được phát hiện kịp thời, dễ dẫn đến thiệt hại cây trồng trên diện rộng.
\end{itemize}

\subsection{Giải pháp hệ thống vườn thông minh (Smart farm)}
Để giải quyết triệt để các vấn đề nêu trên, hệ thống vườn thông minh ứng dụng công nghệ IoT (Internet of Things) được đề xuất như một giải pháp toàn diện. Hệ thống mang lại những ưu điểm vượt trội:
\begin{itemize}
    \item \textbf{Tự động hóa thông minh:} Tích hợp các cảm biến độ ẩm đất để tự động điều tiết bơm nước chính xác theo nhu cầu sinh học của cây, giúp tối ưu hóa tài nguyên nước.
    \item \textbf{Giám sát trực quan:} Dữ liệu được hiển thị trực tiếp trên màn hình LCD và có thể truyền tải qua nền tảng đám mây, giúp người dùng theo dõi trạng thái vườn ở bất cứ đâu.
    \item \textbf{Chế độ điều khiển linh hoạt:} Cho phép chuyển đổi linh hoạt giữa chế độ tự động và điều khiển thủ công từ xa qua điện thoại hoặc máy tính.
\end{itemize}
Việc triển khai hệ thống này không chỉ giúp giảm thiểu công lao động mà còn góp phần hiện đại hóa quy trình sản xuất, hướng tới mô hình nông nghiệp bền vững và thông minh.